\thispagestyle{fancy}
\fancyhead[R]{A.U 2025-2026}
\fancyhead[L]{ chapitre 4: Release 2 : Catalogue et vitrine }
\chapter{Release 2 : Catalogue produit et vitrine personnalisable}
\section*{Introduction}
Ce chapitre présente la deuxième release, qui donne à chaque boutique son contenu et son apparence. Le Sprint 3 met en place un catalogue adapté aux deux secteurs visés, cosmétiques et vêtements, avec des attributs de variantes normalisés. Le Sprint 4 livre trois gabarits de vitrine complets et un éditeur de thème dont chaque réglage s'applique en direct à la vitrine.

%==================================================================
\section{Backlog du Sprint 3}
\Needspace{18\baselineskip}
Le tableau~\ref{tab:backlog-sprint3} présente le \textbf{backlog du Sprint 3}.

\renewcommand{\arraystretch}{1.05}
\setlength{\tabcolsep}{3pt}
\begin{longtable}{|>{\raggedright\arraybackslash}p{4cm}|>{\raggedright\arraybackslash}p{9cm}|>{\centering\arraybackslash}p{1.2cm}|}
\hline
\textbf{User Story} & \textbf{Tâches} & \textbf{Estimation} \\
\hline
\endfirsthead
\hline
\textbf{User Story} & \textbf{Tâches} & \textbf{Estimation} \\
\hline
\endhead
\hline
\multicolumn{3}{r}{\textit{Suite à la page suivante}} \\
\endfoot
\hline
\endlastfoot

\textbf{En tant que commerçant, je souhaite gérer mes produits}
& Développer les entités \texttt{Product} et \texttt{ProductVariant} (prix, stock, statut, images). & 1 \\
\cline{2-3}
& Développer l'API CRUD des produits, avec archivage, désactivation et statistiques. & 1 \\
\cline{2-3}
& Développer l'API de téléversement d'images. & 1 \\
\cline{2-3}
& Développer les interfaces « Produits », « Ajouter un produit » et « Modifier un produit ». & 1 \\
\cline{2-3}
& Adapter le formulaire au secteur de la boutique (champs cosmétiques ou vêtements). & 1 \\
\cline{2-3}
& Tester la fonctionnalité. & 1 \\
\hline

\textbf{En tant que commerçant, je souhaite gérer mes catégories}
& Développer l'API des catégories hiérarchiques (parent, enfants, ordre d'affichage). & 1 \\
\cline{2-3}
& Développer les interfaces de gestion et de réordonnancement des catégories. & 1 \\
\cline{2-3}
& Tester la fonctionnalité. & 1 \\
\hline

\textbf{En tant que commerçant, je souhaite gérer mes collections}
& Développer l'API des collections et l'association des produits. & 1 \\
\cline{2-3}
& Développer les interfaces « Collections » et « Gérer une collection ». & 1 \\
\hline

\textbf{En tant que commerçant, je souhaite choisir les tailles dans une liste}
& Définir les groupes de tailles (lettres, numériques, jeans, enfants, chaussures, taille unique). & 1 \\
\cline{2-3}
& Développer le sélecteur de tailles à choix multiple. & 1 \\
\hline

\textbf{En tant que commerçant, je souhaite choisir des couleurs prédéfinies ou libres}
& Définir la palette de 28 couleurs de base nommées. & 1 \\
\cline{2-3}
& Développer le sélecteur de couleurs avec ajout d'une couleur libre (sélecteur RVB et nom). & 1 \\
\hline

\textbf{En tant que commerçant, je souhaite choisir le tissu et d'autres attributs, et en ajouter}
& Définir la liste des tissus (fibres naturelles, artificielles, synthétiques) et des autres attributs (coupe, saison, entretien). & 1 \\
\cline{2-3}
& Stocker les valeurs ajoutées par la boutique (\texttt{customOptions}) et les fusionner avec les listes par défaut. & 1 \\
\cline{2-3}
& Développer les sélecteurs extensibles et le hook \texttt{useShopOptions}. & 1 \\
\cline{2-3}
& Tester la fonctionnalité. & 1 \\
\hline

\multicolumn{2}{|r|}{\textbf{Total}} & \textbf{19} \\
\hline
\caption{Backlog du Sprint 3}
\label{tab:backlog-sprint3}
\end{longtable}

\section{Diagramme de cas d'utilisation du Sprint 3}
\Needspace{20\baselineskip}
La figure~\ref{fig:cu-sprint3} présente le diagramme de cas d'utilisation du troisième sprint.
\begin{figure}[H]
    \centering
    \img{0.8}{cu_sprint3.png}
    \caption{Diagramme de cas d'utilisation du Sprint 3}
    \label{fig:cu-sprint3}
\end{figure}

\section{Services web du Sprint 3}
Le tableau~\ref{tab:ws-sprint3} présente les services web du Sprint 3, exposés par le microservice \texttt{catalog-service}.
\renewcommand{\arraystretch}{1.15}
\setlength{\tabcolsep}{4pt}
\begin{longtable}{|p{1.5cm}|p{6.8cm}|p{6.2cm}|}
\hline
\textbf{Méthode} & \textbf{URL} & \textbf{Description} \\
\hline
\endfirsthead
\hline
\textbf{Méthode} & \textbf{URL} & \textbf{Description} \\
\hline
\endhead
\multicolumn{3}{|c|}{\textbf{Produits (administration) --- /api/v1/admin/products}} \\
\hline
GET & /api/v1/admin/products & Lister les produits de la boutique. \\
\hline
GET & /api/v1/admin/products/stats & Statistiques du catalogue. \\
\hline
POST & /api/v1/admin/products & Créer un produit et ses variantes. \\
\hline
PUT & /api/v1/admin/products/\{id\} & Modifier un produit. \\
\hline
PATCH & /api/v1/admin/products/\{id\}/archive & Archiver un produit. \\
\hline
PATCH & /api/v1/admin/products/\{id\}/deactivate & Désactiver un produit. \\
\hline
DELETE & /api/v1/admin/products/\{id\} & Supprimer un produit. \\
\hline
POST & /api/v1/admin/upload & Téléverser une image. \\
\hline
\multicolumn{3}{|c|}{\textbf{Catégories et collections (administration)}} \\
\hline
GET, POST & /api/v1/admin/categories & Lister, créer une catégorie. \\
\hline
PUT, DELETE & /api/v1/admin/categories/\{id\} & Modifier, supprimer une catégorie. \\
\hline
PUT & /api/v1/admin/categories/reorder & Réordonner les catégories. \\
\hline
GET, POST & /api/v1/admin/collections & Lister, créer une collection. \\
\hline
PUT, DELETE & /api/v1/admin/collections/\{id\} & Modifier, supprimer une collection. \\
\hline
\multicolumn{3}{|c|}{\textbf{Catalogue public (vitrine)}} \\
\hline
GET & /api/v1/public/products & Lister et filtrer les produits publiés. \\
\hline
GET & /api/v1/public/products/\{slug\} & Fiche d'un produit. \\
\hline
GET & /api/v1/public/products/category/\{categoryId\} & Produits d'une catégorie. \\
\hline
GET & /api/v1/public/products/collection/\{collectionSlug\} & Produits d'une collection. \\
\hline
GET & /api/v1/public/products/by-ids & Produits à partir d'une liste d'identifiants (recommandations). \\
\hline
GET & /api/v1/public/categories/menu & Arborescence du menu de la vitrine. \\
\hline
GET & /api/v1/public/categories/footer & Catégories du pied de page. \\
\hline
GET & /api/v1/public/collections/homepage & Collections mises en avant. \\
\hline
\caption{Services web du Sprint 3}
\label{tab:ws-sprint3}
\end{longtable}

\section{Les cas d'utilisation du Sprint 3}
\subsection{Cas d'utilisation « Ajouter un produit avec ses variantes »}
\subsubsection{Description textuelle}
Le tableau~\ref{tab:uc-produit} présente la description textuelle de ce cas.
\renewcommand{\arraystretch}{1.5}
\begin{longtable}{|p{4cm}|p{11cm}|}
\hline
\textbf{Acteurs} & Commerçant, Collaborateur disposant de la permission « Produits » \\
\hline
\textbf{Objectif} & Publier un produit avec des attributs normalisés, exploitables par les filtres de la vitrine et par le moteur de recommandation. \\
\hline
\textbf{Pré-condition} & L'utilisateur est connecté et sa boutique est active. \\
\hline
\textbf{Scénario principal} &
1. L'utilisateur ouvre « Ajouter un produit » ; le formulaire s'adapte au secteur de la boutique.\newline
2. Il saisit le nom, la description, le prix, la catégorie et téléverse les images.\newline
3. Pour un vêtement, il choisit les tailles dans les groupes proposés, les couleurs dans la palette, et le tissu dans la liste.\newline
4. Le système génère une variante par combinaison taille / couleur, dont l'utilisateur ajuste le stock.\newline
5. L'utilisateur enregistre ; le système valide les données et crée le produit rattaché à sa boutique.\newline
6. Le produit apparaît dans la liste et, s'il est publié, sur la vitrine. \\
\hline
\textbf{Scénario alternatif} &
\textbf{Valeur absente de la liste :} à l'étape 3, l'utilisateur ajoute une couleur libre (sélecteur RVB et nom) ou un nouveau tissu ; la valeur est enregistrée dans les options de la boutique et proposée pour les produits suivants.\newline
\textbf{Données invalides :} à l'étape 5, le système signale les champs erronés (prix négatif, nom manquant) sans perdre la saisie. \\
\hline
\caption{Description textuelle du cas « Ajouter un produit avec ses variantes »}
\label{tab:uc-produit}
\end{longtable}

\subsubsection{Diagramme de séquence système}
\Needspace{18\baselineskip}
\begin{figure}[H]
    \centering
    \img{0.65}{dss_produit.png}
    \caption{Diagramme de séquence système du cas « Ajouter un produit »}
    \label{fig:dss-produit}
\end{figure}

\subsubsection{Diagramme de séquence objet}
\Needspace{20\baselineskip}
\begin{figure}[H]
    \centering
    \img{0.9}{dso_produit.png}
    \caption{Diagramme de séquence objet du cas « Ajouter un produit »}
    \label{fig:dso-produit}
\end{figure}

\section{Réalisation du Sprint 3}
\Needspace{20\baselineskip}
\textbf{Liste des produits :} la figure~\ref{fig:ui-produits} présente la liste des produits avec leurs statistiques, filtres et actions.
\begin{figure}[H]
    \centering
    \img{0.85}{ui_produits.png}
    \caption{Liste des produits de la boutique}
    \label{fig:ui-produits}
\end{figure}

\Needspace{20\baselineskip}
\textbf{Sélecteurs d'attributs :} la figure~\ref{fig:ui-attributs} présente le formulaire d'un vêtement : groupes de tailles, palette de couleurs avec ajout d'une couleur libre, et liste des tissus extensible. Remplacer la saisie libre par ces listes évite les doublons (« bleu marine », « Marine », « navy ») qui fausseraient les filtres et les recommandations.
\begin{figure}[H]
    \centering
    \img{0.85}{ui_attributs.png}
    \caption{Sélecteurs de tailles, couleurs et tissus}
    \label{fig:ui-attributs}
\end{figure}

\Needspace{20\baselineskip}
\textbf{Catégories et collections :} la figure~\ref{fig:ui-categories} présente la gestion des catégories hiérarchiques.
\begin{figure}[H]
    \centering
    \img{0.85}{ui_categories.png}
    \caption{Gestion des catégories}
    \label{fig:ui-categories}
\end{figure}

%==================================================================
\section{Backlog du Sprint 4}
\Needspace{18\baselineskip}
Le tableau~\ref{tab:backlog-sprint4} présente le \textbf{backlog du Sprint 4}, consacré à la vitrine et à sa personnalisation.

\renewcommand{\arraystretch}{1.05}
\setlength{\tabcolsep}{3pt}
\begin{longtable}{|>{\raggedright\arraybackslash}p{4cm}|>{\raggedright\arraybackslash}p{9cm}|>{\centering\arraybackslash}p{1.2cm}|}
\hline
\textbf{User Story} & \textbf{Tâches} & \textbf{Estimation} \\
\hline
\endfirsthead
\hline
\textbf{User Story} & \textbf{Tâches} & \textbf{Estimation} \\
\hline
\endhead
\hline
\multicolumn{3}{r}{\textit{Suite à la page suivante}} \\
\endfoot
\hline
\endlastfoot

\textbf{En tant que commerçant, je souhaite choisir un gabarit Minimal, Bold ou Luxury}
& Concevoir l'architecture des gabarits : en-tête, accueil et pied de page par gabarit, briques partagées (navigation, données d'accueil, catégories du pied de page). & 1 \\
\cline{2-3}
& Développer le gabarit Minimal (épuré, typographie fine, grands espaces). & 1 \\
\cline{2-3}
& Développer le gabarit Bold (contrastes forts, en-tête transparent sur bannière plein écran). & 1 \\
\cline{2-3}
& Développer le gabarit Luxury (typographie à empattements, tons sourds, mise en page éditoriale). & 1 \\
\cline{2-3}
& Faire correspondre les anciens thèmes aux trois gabarits et permettre la prévisualisation (\texttt{?preview=}). & 1 \\
\cline{2-3}
& Supprimer la section codée en dur propre à NaturEssence (atelier « DIY »). & 1 \\
\hline

\textbf{En tant que visiteur, je souhaite parcourir la vitrine, ses menus et ses catégories}
& Développer les menus déroulants alimentés par l'arborescence des catégories. & 1 \\
\cline{2-3}
& Développer la page de catégorie avec filtres (prix, taille, couleur). & 1 \\
\cline{2-3}
& Afficher un message lorsque la boutique est en attente, fermée ou suspendue. & 1 \\
\hline

\textbf{En tant que visiteur, je souhaite consulter la fiche d'un produit}
& Développer la fiche produit (galerie, variantes, stock, avis). & 1 \\
\cline{2-3}
& Adapter la carte produit à chaque gabarit, avec choix rapide de la taille. & 1 \\
\hline

\textbf{En tant que commerçant, je souhaite choisir la palette de couleurs de ma boutique}
& Exposer la palette sous forme de variables CSS (\texttt{--rgb-primary}, \texttt{--rgb-accent}, \dots) consommées par Tailwind. & 1 \\
\cline{2-3}
& Développer le choix de palette dans « Votre boutique » (couleurs prédéfinies et personnalisées). & 1 \\
\hline

\textbf{En tant que commerçant, je souhaite personnaliser boutons, survols, navigation, textes et pied de page}
& Définir le modèle de thème (14 variables : annonce, navigation, survol, titres, boutons, prix, badges, pied de page). & 1 \\
\cline{2-3}
& Développer l'éditeur de thème avec aperçu et réinitialisation. & 1 \\
\cline{2-3}
& Appliquer le thème sur la vitrine via des variables \texttt{--t-*} et des classes d'ancrage (\texttt{t-nav}, \texttt{t-btn}, \texttt{t-price}\dots). & 1 \\
\cline{2-3}
& Valider le thème côté serveur (couleurs \#RRGGBB, taille et type du logo). & 1 \\
\cline{2-3}
& Tester sur les trois gabarits. & 1 \\
\hline

\multicolumn{2}{|r|}{\textbf{Total}} & \textbf{18} \\
\hline
\caption{Backlog du Sprint 4}
\label{tab:backlog-sprint4}
\end{longtable}

\section{Diagramme de cas d'utilisation du Sprint 4}
\Needspace{20\baselineskip}
\begin{figure}[H]
    \centering
    \img{0.8}{cu_sprint4.png}
    \caption{Diagramme de cas d'utilisation du Sprint 4}
    \label{fig:cu-sprint4}
\end{figure}

\section{Services web du Sprint 4}
Le tableau~\ref{tab:ws-sprint4} présente les services web du Sprint 4.
\renewcommand{\arraystretch}{1.15}
\setlength{\tabcolsep}{4pt}
\begin{longtable}{|p{1.5cm}|p{6.8cm}|p{6.2cm}|}
\hline
\textbf{Méthode} & \textbf{URL} & \textbf{Description} \\
\hline
\endfirsthead
\hline
\textbf{Méthode} & \textbf{URL} & \textbf{Description} \\
\hline
\endhead
\multicolumn{3}{|c|}{\textbf{Boutique et thème (auth-service)}} \\
\hline
PATCH & /api/v1/auth/my-shop & Enregistrer le gabarit, la palette, le thème et le logo. \\
\hline
GET & /api/v1/public/shops/\{slug\} & Charger le gabarit, la palette et le thème d'une boutique. \\
\hline
\multicolumn{3}{|c|}{\textbf{Apparence et profil de la boutique (marketing-service)}} \\
\hline
GET & /api/v1/admin/appearance/\{scope\} & Consulter les réglages d'apparence d'une zone (accueil, en-tête, pied de page). \\
\hline
PUT & /api/v1/admin/appearance/\{scope\} & Modifier les réglages d'apparence. \\
\hline
POST & /api/v1/admin/appearance/\{scope\}/reset & Réinitialiser les réglages. \\
\hline
GET & /api/v1/public/appearance/\{scope\} & Réglages d'apparence publics. \\
\hline
GET, PUT & /api/v1/admin/store & Consulter, modifier le profil de la boutique (contact, réseaux sociaux). \\
\hline
GET & /api/v1/public/store & Profil public de la boutique. \\
\hline
\caption{Services web du Sprint 4}
\label{tab:ws-sprint4}
\end{longtable}

\section{Les cas d'utilisation du Sprint 4}
\subsection{Cas d'utilisation « Personnaliser la vitrine »}
\subsubsection{Description textuelle}
Le tableau~\ref{tab:uc-theme} présente la description textuelle de ce cas.
\renewcommand{\arraystretch}{1.5}
\begin{longtable}{|p{4cm}|p{11cm}|}
\hline
\textbf{Acteurs} & Commerçant \\
\hline
\textbf{Objectif} & Donner à la vitrine l'apparence de la marque sans écrire de code. \\
\hline
\textbf{Pré-condition} & Le commerçant est connecté et possède une boutique. \\
\hline
\textbf{Scénario principal} &
1. Le commerçant ouvre « Votre boutique ».\newline
2. Il choisit un gabarit (Minimal, Bold ou Luxury) et peut le prévisualiser sans l'enregistrer.\newline
3. Il choisit sa palette : couleur de marque, accent, fond, texte.\newline
4. Dans l'éditeur de thème, il règle le bandeau d'annonce, la navigation et son survol, les boutons et leur survol, les prix, les badges et le pied de page.\newline
5. Il enregistre ; le système valide chaque couleur (format \#RRGGBB) et stocke le thème.\newline
6. À l'ouverture suivante, la vitrine charge le thème et l'applique sous forme de variables CSS ; seuls les réglages renseignés remplacent le style par défaut du gabarit. \\
\hline
\textbf{Scénario alternatif} &
\textbf{Couleur invalide :} à l'étape 5, le système refuse le thème avec « Le thème ne peut contenir que des couleurs \#RRGGBB ».\newline
\textbf{Retour au style du gabarit :} le commerçant réinitialise un réglage ; la variable correspondante est retirée et le style d'origine du gabarit s'applique de nouveau. \\
\hline
\caption{Description textuelle du cas « Personnaliser la vitrine »}
\label{tab:uc-theme}
\end{longtable}

\subsubsection{Diagramme de séquence système}
\Needspace{18\baselineskip}
\begin{figure}[H]
    \centering
    \img{0.65}{dss_theme.png}
    \caption{Diagramme de séquence système du cas « Personnaliser la vitrine »}
    \label{fig:dss-theme}
\end{figure}

\subsubsection{Diagramme de séquence objet}
\Needspace{20\baselineskip}
\begin{figure}[H]
    \centering
    \img{0.9}{dso_theme.png}
    \caption{Diagramme de séquence objet du cas « Personnaliser la vitrine »}
    \label{fig:dso-theme}
\end{figure}

\section{Réalisation du Sprint 4}
\Needspace{20\baselineskip}
\textbf{Les trois gabarits :} les figures~\ref{fig:ui-minimal}, \ref{fig:ui-bold} et~\ref{fig:ui-luxury} présentent la même boutique affichée avec chacun des trois gabarits. Le contenu (produits, catégories, bannières) est identique ; seuls la structure et le style changent.
\begin{figure}[H]
    \centering
    \img{0.85}{ui_vitrine_minimal.png}
    \caption{Vitrine avec le gabarit Minimal}
    \label{fig:ui-minimal}
\end{figure}
\begin{figure}[H]
    \centering
    \img{0.85}{ui_vitrine_bold.png}
    \caption{Vitrine avec le gabarit Bold}
    \label{fig:ui-bold}
\end{figure}
\begin{figure}[H]
    \centering
    \img{0.85}{ui_vitrine_luxury.png}
    \caption{Vitrine avec le gabarit Luxury}
    \label{fig:ui-luxury}
\end{figure}

\Needspace{20\baselineskip}
\textbf{Éditeur de thème :} la figure~\ref{fig:ui-theme} présente l'éditeur de thème, organisé en groupes (couleurs de base, bandeau d'annonce, navigation, boutons, prix et badges, pied de page).
\begin{figure}[H]
    \centering
    \img{0.85}{ui_editeur_theme.png}
    \caption{Éditeur de thème de la vitrine}
    \label{fig:ui-theme}
\end{figure}

\Needspace{20\baselineskip}
\textbf{Fiche produit :} la figure~\ref{fig:ui-fiche} présente la fiche d'un vêtement, avec le choix de la taille et de la couleur.
\begin{figure}[H]
    \centering
    \img{0.85}{ui_fiche_produit.png}
    \caption{Fiche produit sur la vitrine}
    \label{fig:ui-fiche}
\end{figure}

\subsection*{Conclusion}
Cette release a doté chaque boutique d'un catalogue structuré et d'une vitrine qui lui ressemble. Les attributs normalisés ne servent pas qu'à l'affichage : ils alimentent aussi la partie « contenu » du moteur de recommandation présenté au chapitre 8. La release suivante porte sur le parcours d'achat.
